\section{Localization, occupation, and isolation}\label{sec:localization}

We now obtain simultaneous probability bounds for specified collections
of cells. The important scale is $s\asymp K\rho^{-1/2}$: a cell contains
$m=Ns^3/V\asymp\rho s^3$ particles on average, and $m$ tends to infinity
as the density decreases. All cutoffs below are real and Lipschitz;
their gradients and the resulting form identities are understood
almost everywhere.

The estimates use the same variational comparison. Spatial localization
gives a lower bound for the free-energy cost of a tilted state, while
diagonal tilting gives an upper bound for that cost in terms of a gradient
error and the tilt's normalization. Comparing the two bounds controls
the derivative of a logarithmic exponential moment. Positive density
tilts control excess occupation, negative tilts control shortages, and
a tilt of the close-particle count will give isolation.

\begin{lemma}[Two-color localization and diagonal tilting]\label{lem:localization-tilt}
Let $p^2+q^2=1$, and suppose multiplication by $p$ maps the one-particle
form domain into $H^1_0(\Omega)$ for an open region $\Omega$.
For a finite-energy symmetric trace-one $J$-particle state $\sigma$, let $w_j$ be
the distribution of the number of particles in the $p$ color, and
let $\sigma_{\mathrm{in},j},\sigma_{\mathrm{out},j}$ be the two
conditional marginal states. Then
\begin{equation}
 \mathcal F^{(-W)}(\sigma)+\operatorname{Tr}
 n_{|\nabla p|^2+|\nabla q|^2}\sigma
 \geq\sum_{j=0}^Jw_j\bigl[
 \mathcal F^{(-W)}(\sigma_{\mathrm{in},j})+
 \mathcal F^{(-W)}(\sigma_{\mathrm{out},j})\bigr]-TH(w).
 \tag{7}\label{eq:7}
\end{equation}
Here $H(w)=-\sum_jw_j\log w_j$, with $0\log0=0$, is the Shannon
entropy, and $H(w)\leq C+\sum_jjw_j$ and
$\sum_jjw_j=\operatorname{Tr}n_{p^2}\sigma$.

Let $\Gamma=Z_N^{-1}e^{-\beta H_{N,L}}$, let $P$ be its position
law, and let $G$ be a bounded symmetric Lipschitz real function of
positions. Put $g=e^G$, $M_G=\mathbb E_Pe^{2G}$ and
$\sigma=g\Gamma g/M_G$. For expectation under its position law
write $\mathbb E_g$. Then
\begin{equation}
 \mathcal F(\sigma)-F_N\leq
 \mathbb E_g\sum_i|\nabla_iG|^2+
 T\bigl(2\mathbb E_gG-\log M_G\bigr).
 \tag{8}\label{eq:8}
\end{equation}
\end{lemma}

\begin{proof}
The geometric localization and entropy split behind (7) also appear in
\cite[Section~3, Lemma~3.1]{DeuchertSeiringerYngvason2019}.
We give the decomposition into fixed color-number sectors explicitly.
The one-particle map $f\mapsto(pf,qf)$ is an isometry into the direct
sum of the two one-particle spaces. Its symmetric tensor power is
therefore an isometry. Dephase the resulting state with respect to
the first-color number. The $j$-sector is unitarily identified with
$\operatorname{Sym}^jL^2(\Omega)\otimes
\operatorname{Sym}^{J-j}L^2(\Lambda_L)$: normalized symmetrization
over color assignments gives this unitary, because the assignments
are orthogonal. In particular there is no additional multiplicity
or binomial entropy term. The product rule and
$p\nabla p+q\nabla q=0$ give exactly the kinetic error in (7).
Retain the same-color pair interactions and drop the other ones;
nonnegativity makes this an energy decrease. One-body potentials
are preserved in sum. Isometry preserves entropy, dephasing cannot
decrease it, and subadditivity in each sector bounds it by
$H(w)+\sum_jw_j[S(\sigma_{\mathrm{in},j})+
S(\sigma_{\mathrm{out},j})]$. This proves (7).
Comparison of $w$ with the geometric law
$(1-e^{-1})e^{-j}$ gives $H(w)\leq-\log(1-e^{-1})+\mathbb E j$.
Conditional on a labeled position configuration, the color choices
are independent, with probabilities $p(x_i)^2$ and $q(x_i)^2$.
This proves the number formula and also the diagonal interpretation
of every color count used below. Finite-energy approximation and
Lemma~\ref{lem:kinetic-entropy} justify all the entropy operations.

For (8), define
$\tau=\Gamma^{1/2}g^2\Gamma^{1/2}/M_G$.
The nonzero spectra of $\tau$ and $\sigma$ agree: they are the
spectra of $B^*B$ and $BB^*$ for $B=g\Gamma^{1/2}/\sqrt{M_G}$.
If $H\psi=E\psi$, testing the eigenvalue equation with $g^2\psi$
gives the form identity
\[
 \langle g\psi,Hg\psi\rangle
 =E\langle\psi,g^2\psi\rangle+
 \langle\psi,\sum_i|\nabla_i g|^2\psi\rangle.
\]
After summing in an eigenbasis,
$\operatorname{Tr}H\sigma=\operatorname{Tr}H\tau+
\mathbb E_g\sum_i|\nabla_iG|^2$. Both energies and entropies are
finite, since $g$ and its gradients are bounded at fixed $N,L$.

To compare relative entropies, approximate $g^2$ uniformly by positive
simple functions on finite partitions into symmetric measurable
sets $E$. Set $p_E=\operatorname{Tr}\Gamma\mathbf1_E$ and, for
$p_E>0$, $\tau_E=\Gamma^{1/2}\mathbf1_E\Gamma^{1/2}/p_E$.
The unmodified and tilted states are mixtures of the same $\tau_E$,
with weights $p_E$ and $p_Eg_E^2/M_G$. Monotonicity of relative
entropy under discarding the classical label (a partial trace;
see \cite[Corollary~(3.2) and Section~IV]{LiebRuskai1973Entropy}) therefore bounds
$D(\tau\Vert\Gamma)$ for this approximation by the classical
relative entropy of these weights. Uniform approximation gives
trace-norm convergence of the states and convergence of the
classical entropies, since $g$ is bounded above and away from zero.
Lower semicontinuity then gives
$D(\tau\Vert\Gamma)\leq2\mathbb E_gG-\log M_G$.
Finally $\mathcal F(\tau)-F_N=TD(\tau\Vert\Gamma)$ proves (8).
\end{proof}

For a nonempty union $U$ of coarse cells, set
$\phi_U(x)=e^{-d(x,U)/s}$, where distance is the Euclidean torus
distance. A neighborhood of radius $hs$ around $U$ is covered by
at most $C(h+2)^3$ times as many cells as $U$.

\begin{lemma}[Localized density bounds]\label{lem:density-tilts}
Let $D_0=H_0\rho$, where $H_0$ is a sufficiently large fixed number.
For $J\leq N$, $0\leq W\leq\rho$, and every finite-energy
symmetric trace-one $J$-particle state,
\begin{equation}
 \mathcal F^{(-W)}(\sigma)-F_J^{(-W)}
 \geq(c_vD_0-C_v\rho-Cs^{-2})\mathbb E_\sigma n_{\phi_U}
       -C_vD_0^2|U|.\tag{9}\label{eq:9}
\end{equation}
In the corresponding equilibrium state,
\begin{equation}
 \mathbb E n_U\leq C_v\rho|U|.\tag{10}\label{eq:10}
\end{equation}
For the unmodified $N$-particle position law,
\begin{equation}
 \log\mathbb E_Pe^{2\lambda n_{\phi_U}}
 \leq C_v\lambda\rho|U|\qquad(0\leq\lambda\leq1).
 \tag{11}\label{eq:11}
\end{equation}
In particular $\mathbb E_P n_Q^2\leq C m^2$ for every cell $Q$.
The proof of (9) can reserve a fixed positive fraction of the
inside kinetic and pair energy for a further nonnegative lower bound.
\end{lemma}

\begin{proof}
For $h\geq0$ take $p_h=1$ at distance at most $hs$ from $U$,
and interpolate by a cosine to zero between distances $hs$ and
$(h+1)s$. Set $q_h=(1-p_h^2)^{1/2}$ using the corresponding sine.
Both gradients vanish off this transition and
$|\nabla p_h|^2+|\nabla q_h|^2\leq Cs^{-2}$ there. Choose
an open support region $\Omega_h$ inside the $(h+2)s$
neighborhood; if that neighborhood covers the torus, use the torus.
Its volume is at most $C(h+4)^3|U|$.

Apply (7). Inside, partition the torus into cubes of side in
$[D_0^{-1/3},2D_0^{-1/3}]$. For small density this side is between
$b_0$ and $s$. The cubes meeting $\Omega_h$ number at most
$CD_0(h+4)^3|U|$. Applying (2) to part of the inside energy of
zero-extended Dirichlet states gives
$c_vD_0j-CD_0^2(h+4)^3|U|$.
Another fixed fraction of kinetic energy gives (3). An additional
fixed fraction may remain unused. The external potential costs at
most $\rho j$. Outside, (6) gives
$F_{J-j}^{(-W)}\geq F_J^{(-W)}-C_v\rho j$.
The classical entropy of the color number costs at most $T(C+j)$.
Since $T\leq\rho^2$, all coefficients of $j$ other than
$c_vD_0$ are bounded by $C_v\rho$. All volume errors, including
$CT[1+|\Omega_h|(V^{-1}+T^{3/2})]$, are at most
$C_vD_0^2(h+4)^3|U|$ after reducing the density: here
$|U|\geq s^3$ and $V\geq s^3$.

Average in $h$ with probability density $e^{-h}\,dh$. At each $x$,
$\int_0^\infty p_h(x)^2e^{-h}\,dh\geq\phi_U(x)$.
The averaged gradient cost is at most $Cs^{-2}\phi_U(x)$:
the transition can occur only for
$h\in[\max(0,d(x,U)/s-1),d(x,U)/s]$.
The moment of $(h+4)^3$ is finite. Taking $H_0$ large makes
$c_vD_0-C_v\rho$ positive and yields (9).
In equilibrium the left side is zero; because
$n_U\leq n_{\phi_U}$ this proves (10).

For (11), take $G=\lambda n_{\phi_U}$ in (8).
The distance function is one-Lipschitz, so
$\sum_i|\nabla_iG|^2\leq\lambda^2s^{-2}n_{\phi_U}$.
Also $G\geq0$, hence $\log M_G\geq0$ and the entropy term is
at most $2T\lambda\mathbb E_g n_{\phi_U}$.
Comparison with (9), increasing the fixed $H_0$ if necessary,
absorbs both terms and gives
$\mathbb E_g n_{\phi_U}\leq C_v\rho|U|$ uniformly in
$0\leq\lambda\leq1$. Integrate
$\partial_\lambda\log M_G=2\mathbb E_g n_{\phi_U}$
from zero to obtain (11).
For a single cell, exponential Markov gives
$P(n_Q\geq t)\leq\min(1,e^{Cm-2t})$.
Integration of $2tP(n_Q>t)$ gives $\mathbb E n_Q^2\leq C(m+1)^2$,
which is at most $Cm^2$ since $m\to\infty$ uniformly with dilution.
\end{proof}

The mean bound alone does not guarantee that a cell supplies a particle
to a spatial route. We next control simultaneous shortages and excesses
in any specified family of cells. The lower tail uses the free-energy
gain from placing extra particles in an underoccupied region; the upper
tail follows from the exponential density estimate.

\begin{lemma}[Simultaneous lower and upper occupation tails]\label{lem:occupation-tails}
There are fixed $0<b_-<1/4$ and $b_+>4$ such that for every union
$U$ of $u\geq1$ specified cells, each of the events
$\{n_Q<2b_-m\text{ for every }Q\subset U\}$ and
$\{n_Q>b_+m/2\text{ for every }Q\subset U\}$ has probability at most
\begin{equation}
 C e^{-cmu}.\tag{12}\label{eq:12}
\end{equation}
The constants and the required large fixed $K$ depend only on $v$.
\end{lemma}

\begin{proof}
In each cell of $U$ choose $p=1$ on the central cube of half-side
$s/4$, vanishing outside the central cube of half-side $s/3$.
Cosine--sine interpolation gives squared localization cost at most
$Cs^{-2}$ on these supports. Choose a squared smooth cutoff
$0\leq\phi\leq1$, supported in the cell interiors, equal to one
on these supports, with $|\nabla\phi|^2\leq C\phi/s^2$.
The union of cell interiors is the Dirichlet region for $p$.
Take one normalized smooth bump in each central cube of half-side
$s/8$ and sum the bumps with equal amplitude to get a normalized
orbital $\varphi$. Then
$\int|\nabla\varphi|^2\leq Cs^{-2}$ and
$\|\varphi\|_\infty^2\leq C/|U|$.
For small density, its support has distance greater than $R_v$
from $\operatorname{supp}q$.

The conditional outside state is supported in
$(\operatorname{supp}q)^{N-j}$, directly from the coloring map.
In a color sector with $j$ inside particles, add $k$ particles in
$\varphi$ to the outside state. The map on outside wave functions is
\[
 \psi\longmapsto
 \binom{N-j+k}{k}^{1/2}P_{\mathrm{sym}}
 (\psi\otimes\varphi^{\otimes k}).
\]
Distinct assignments of the new labels have disjoint supports, so
this is an isometry. It preserves entropy and has additive kinetic
energy; the separated supports also eliminate all old--new
interactions. A new--new pair has expectation at most
$\|v\|_1\|\varphi\|_\infty^2$. Therefore
\[
 \mathcal F(\sigma_{\mathrm{out},j})
 \geq F_{N-j+k}-Ck/s^2-C_vk^2/|U|.
\]
Choose a sufficiently small fixed $\eta>0$ and put
$k=\lfloor\eta\rho|U|\rfloor$.
Since $\rho|U|\geq\rho s^3\to\infty$, small density ensures
$k\geq\eta\rho|U|/2$, and $N+k\leq4\rho V$.
By (5), followed by (6) for $j$ backward steps,
$F_{N-j+k}\geq F_N+c_v\rho k-C_v\rho j$.
Use (3) and positivity inside. Choose $\eta$ small enough to absorb
$C_vk^2/|U|$ and then $K$ large enough to absorb $Ck/s^2$ into a
fixed fraction of $c_v\rho k$. The remaining entropy errors are
small compared with $\rho^2|U|$ for small density.
Applying (7) and using $p^2\leq\phi$ gives, for every state to
which the argument is applied,
\[
 \mathcal F(\sigma)-F_N
 \geq c_v'\rho^2|U|-C_v\rho\mathbb E_\sigma n_\phi.
\]
Here the localization term has been absorbed in the latter error,
since $s^{-2}\leq\rho/K^2$.

Now set $G=-\lambda n_\phi$, $0\leq\lambda\leq1$, in (8).
Its gradient contribution is at most
$C\lambda^2s^{-2}\mathbb E_gn_\phi$; its entropy contribution is
at most $-T\log M_G$, because $G\leq0$.
Thus, after changing constants,
$C_v\rho\mathbb E_g n_\phi-T\log M_G\geq c_v'\rho^2|U|$.
Put $f(\lambda)=-\log M_G$. It is nondecreasing and
$f'(\lambda)=2\mathbb E_gn_\phi$. While
$f(\lambda)\leq a\rho|U|$, one has
$Tf(\lambda)\leq a\rho^3|U|$, so for sufficiently small density
$f'(\lambda)\geq c_v''\rho|U|$. Choose fixed
$0<a<c_v''$. Starting at $f(0)=0$, this differential inequality
forces $f(1)\geq a\rho|U|$. We have proved
$\mathbb E e^{-2n_\phi}\leq e^{-a\rho|U|}$.

On the lower-count event, $n_\phi\leq2b_-mu$, so its probability
is at most $e^{4b_-mu-a\rho|U|}$. Since
$mu/2\leq\rho|U|\leq2mu$, sufficiently small $b_-$ gives
$e^{-cmu}$. On the upper-count event,
$n_{\phi_U}\geq b_+mu/2$; (11) at $\lambda=1$ bounds its
probability by $e^{C_v\rho|U|-b_+mu}$. A sufficiently large
$b_+$ proves the other bound in (12).
\end{proof}

Occupation control supplies many particles per cell. The bridge
comparison will additionally require well-separated endpoints. The
next estimate shows that a prescribed region rarely contains many
particles with a neighbor at the shorter scale $r$.

\begin{lemma}[Isolation at a shorter scale]\label{lem:isolation}
Fix $0<\zeta\leq1$ and set $r=\zeta\rho^{-1/3}$ and
$\lambda_0=r\sqrt\rho$. A particle is nonisolated if another
particle is within torus distance $r$. For any nonempty union $S$
of coarse cells and any $k'\geq0$,
\begin{equation}
 P(\#\{\text{nonisolated particles in }S\}\geq k')
 \leq\exp\{C_v\lambda_0\zeta^3\rho|S|-\lambda_0 k'\}.
 \tag{13}\label{eq:13}
\end{equation}
The constant $C_v$ can be chosen independently of decreasing
$\zeta$; the density threshold is allowed to depend on $\zeta$.
\end{lemma}

\begin{proof}
Partition the torus into cubes of side in $[r,2r]$.
For each cube meeting $S$ choose a cutoff $0\leq\chi\leq1$ that is one on
the concentric triple cube, vanishes outside the concentric
quintuple cube, and satisfies $|\nabla\chi|\leq C/r$.
Let $Q=\sum_\chi(n_\chi-1)_+$. A nonisolated particle in a
fine cube has a partner in that cube's triple cube. If there are
$a$ nonisolated particles in the original cube, its cutoff has
$n_\chi\geq\max(a,2)$, and hence
$a\leq2(n_\chi-1)_+$. Summing gives the deterministic bound
$\#\{\text{nonisolated in }S\}\leq2Q$.
Bounded overlap of the quintuple cubes gives
$\sum_i|\nabla_i Q|^2\leq Cr^{-2}n_{U_0}$,
where $U_0$ is a coarse-cell cover of their union with
$|U_0|\leq C|S|$. The cubes admit a coloring with a fixed number
of colors, disjoint within each color. This can be obtained from
the bounded-degree intersection graph and works also on the torus.

Repeat the proof of (9) for $U_0$, reserving a fixed fraction of
inside energy for (2) on these quintuple cubes, divided among the
colors. In every localization mask $p_h=1$ on $U_0$, so those
particles are certainly assigned to the inside color. Also
$(n_\chi-1)_+\leq n_{\operatorname{supp}\chi}
\mathbf1_{\{n_{\operatorname{supp}\chi}\geq2\}}$.
The reserved energy therefore contributes
$c_vr^{-3}\mathbb E_\sigma Q$. The small-density conditions ensure
that the fine cubes exceed $b_0$; all overlap constants are
independent of $\zeta$.

Apply this augmented lower bound to $G=\lambda Q$ for
$0\leq\lambda\leq\lambda_0$. Its gradient cost in (8) is
at most $C\rho\mathbb E_g n_{U_0}$, absorbed by the density term
of (9) after fixing $H_0$ sufficiently large. The entropy cost is
at most $2T\lambda\mathbb E_gQ$, absorbed by
$c_vr^{-3}\mathbb E_gQ$ for small density. It follows that
$\mathbb E_gQ\leq C_vr^3\rho^2|S|
=C_v\zeta^3\rho|S|$.
Integration of $\partial_\lambda\log\mathbb E e^{2\lambda Q}
=2\mathbb E_gQ$ yields
$\log\mathbb E e^{2\lambda_0Q}
\leq C_v\lambda_0\zeta^3\rho|S|$.
Since $2Q\geq\#\{\text{nonisolated in }S\}$, exponential Markov
proves (13).
\end{proof}
